\chapter{Identificación de leyes y regulaciones}

\section{Normativa relevante}
Odín trabaja con información personal, documentos privados, fotografías, sensores domésticos y cámaras. Por ello, aunque sea un proyecto académico y de uso personal, conviene identificar el marco regulatorio aplicable:
\begin{itemize}
  \item el Reglamento General de Protección de Datos (RGPD), aplicable al tratamiento de datos personales en la Unión Europea;
  \item la Ley Orgánica 3/2018 de Protección de Datos Personales y garantía de los derechos digitales (LOPDGDD), que adapta y completa el RGPD en España;
  \item las orientaciones de la Agencia Española de Protección de Datos sobre videovigilancia, relevantes cuando las cámaras puedan captar a terceros \citep{gdprOfficial,lopdgddOfficial,aepdVideovigilancia}.
\end{itemize}

\section{Aplicación al proyecto}
La filosofía \textit{local-first} favorece el cumplimiento por diseño: los datos se conservan en infraestructura propia, se reduce la transferencia a terceros y se limita la exposición pública. Sin embargo, autoalojar no elimina responsabilidades. El proyecto debe seguir prestando atención a:
\begin{itemize}
  \item minimización de datos almacenados;
  \item control de accesos y protección de credenciales;
  \item conservación limitada de imágenes o eventos cuando intervengan terceros;
  \item información y señalización adecuadas si la videovigilancia sale del ámbito puramente doméstico;
  \item copias de seguridad protegidas y posibilidad de supresión cuando proceda.
\end{itemize}

\section{Cumplimiento y decisiones técnicas}
Varias decisiones de Odín se alinean con estas obligaciones:
\begin{itemize}
  \item evitar publicar secretos en repositorios;
  \item preferir VPN frente a exposición pública innecesaria;
  \item separar credenciales en ficheros de entorno;
  \item documentar configuraciones saneadas;
  \item mantener cámaras y servicios bajo control del usuario;
  \item tratar las funcionalidades de vigilancia como un caso sensible, no como un mero experimento técnico.
\end{itemize}

\section{Límites}
Este informe no sustituye un análisis jurídico profesional. Su objetivo es demostrar que el proyecto identifica las normas relevantes y adopta decisiones técnicas coherentes con ellas. Si el sistema se extendiera a usuarios externos, espacios compartidos o servicios ofrecidos a terceros, sería necesario revisar con mayor profundidad bases legitimadoras, deberes de información, contratos y responsabilidades adicionales.
